% --- Archon physics formalization source begin ---
% archon:physics
% archon:covers PhyXMiniProblems/problem_phyx_mini_0666.lean
% archon:source-report reports/phyx_mini/problem_phyx_mini_0666.source.json

\chapter{Physics problem phyx\_mini\_0666}
\label{ch:PhyXMiniProblems_problem_phyx_mini_0666}

\paragraph{Problem source.}
\#\# Physical scenario

Figure shows a frustum of a cone.

\#\# Question

Calculate the area of the curved surface.

\#\# Answer choices

- A: A: \textbackslash{}[ \textbackslash{}3pi(r\_1 + r\_2)[h\textasciicircum{}\{2\} + (r\_2 - r\_1)\textasciicircum{}\{2\}]\textasciicircum{}\{1/2\} \textbackslash{}]

- B: B: \textbackslash{}[ \textbackslash{}2pi(r\_1 + r\_2)[h\textasciicircum{}\{2\} + (r\_2 - r\_1)\textasciicircum{}\{2\}]\textasciicircum{}\{1/2\} \textbackslash{}]

- C: C: \textbackslash{}[ \textbackslash{}pi(r\_1 + r\_2)[h\textasciicircum{}\{2\} + (r\_2 - r\_1)\textasciicircum{}\{2\}]\textasciicircum{}\{1/2\} \textbackslash{}]

- D: D: \textbackslash{}[ \textbackslash{}4pi(r\_1 + r\_2)[h\textasciicircum{}\{2\} + (r\_2 - r\_1)\textasciicircum{}\{2\}]\textasciicircum{}\{1/2\} \textbackslash{}]

\#\# Figure caption (auxiliary; use the image as primary evidence)

The image is a diagram of a truncated cone, also known as a frustum. The cone is three-dimensional and has two parallel circular bases. The top base has a radius labeled \textbackslash{}( r\_1 \textbackslash{}), and the bottom base has a radius labeled \textbackslash{}( r\_2 \textbackslash{}). The height of the frustum, perpendicular to the bases, is denoted by \textbackslash{}( h \textbackslash{}).

Dashed lines indicate the radial lines \textbackslash{}( r\_1 \textbackslash{}) and \textbackslash{}( r\_2 \textbackslash{}), along with the height \textbackslash{}( h \textbackslash{}). The overall appearance suggests the frustum is transparent or semi-transparent to emphasize these dimensions.

\#\# Dataset metadata

Statics; Spatial Relation Reasoning

\paragraph{Recorded answer/context.}
C: C: \textbackslash{}[ \textbackslash{}pi(r\_1 + r\_2)[h\textasciicircum{}\{2\} + (r\_2 - r\_1)\textasciicircum{}\{2\}]\textasciicircum{}\{1/2\} \textbackslash{}]

\paragraph{Figure/image path.}
/root/proposal\_for\_physic/science-mango/phyx\_mini\_run/phyx\_data/test\_image/666.png

\paragraph{Formalization target.}
create a compiling Lean file with sorry bodies at `PhyXMiniProblems/problem\_phyx\_mini\_0666.lean`. The Lean declarations must preserve the physical quantities, dimensions or dimensional roles, figure labels, governing-law hypotheses, and final relation expressed by this problem.
Use Mathlib/Physlib names found through LeanExplore where available. If a physics API is missing, introduce faithful local abstractions rather than scalar placeholder aliases.

\begin{theorem}[Physics formalization target]
\label{thm:physics:phyx_mini_0666:target}
\lean{PhyXMiniProblems.ProblemPhyXMini0666.problem_phyx_mini_0666}
\uses{def:physics:phyx-mini-0666:phyxminiproblems-problemphyxmini0666-lengthinmeters, def:physics:phyx-mini-0666:phyxminiproblems-problemphyxmini0666-areainsquaremeters, def:physics:phyx-mini-0666:phyxminiproblems-problemphyxmini0666-rightcircularconicalfrustum, def:physics:phyx-mini-0666:phyxminiproblems-problemphyxmini0666-matchessuppliedfrustumfigure, def:physics:phyx-mini-0666:phyxminiproblems-problemphyxmini0666-hasvalidfrustumparameters, def:physics:phyx-mini-0666:phyxminiproblems-problemphyxmini0666-satisfiesrightfrustummeridiangeometry, def:physics:phyx-mini-0666:phyxminiproblems-problemphyxmini0666-satisfiesfrustumcurvedsurfacearealaw, def:physics:phyx-mini-0666:phyxminiproblems-problemphyxmini0666-answerchoice, def:physics:phyx-mini-0666:phyxminiproblems-problemphyxmini0666-displayedanswervalue}
The assigned autoformalize agent should translate this physics problem into Lean declarations in the covered file, with theorem and lemma proof bodies written as `by sorry`.
\end{theorem}
\begin{proof}
This is an autoformalization task, not a proof task. Produce statements that can later be proved by the physics prover without weakening the physical model.
\end{proof}

% --- Archon declaration topology begin ---
\section{Declaration topology}
\begin{definition}[Length Quantity]
  \label{def:physics:phyx-mini-0666:phyxminiproblems-problemphyxmini0666-lengthquantity}
  \lean{PhyXMiniProblems.ProblemPhyXMini0666.LengthQuantity}
  A nonnegative, unit-independent physical length
\end{definition}
\begin{proof}
  This is the declared model definition of Length Quantity.
\end{proof}
\begin{definition}[length In Meters]
  \label{def:physics:phyx-mini-0666:phyxminiproblems-problemphyxmini0666-lengthinmeters}
  \lean{PhyXMiniProblems.ProblemPhyXMini0666.lengthInMeters}
  \uses{def:physics:phyx-mini-0666:phyxminiproblems-problemphyxmini0666-lengthquantity}
  Metre readout of a physical length
\end{definition}
\begin{proof}
  This is the declared model definition of length In Meters.
\end{proof}
\begin{definition}[area In Square Meters]
  \label{def:physics:phyx-mini-0666:phyxminiproblems-problemphyxmini0666-areainsquaremeters}
  \lean{PhyXMiniProblems.ProblemPhyXMini0666.areaInSquareMeters}
  Square-metre readout of a physical area
\end{definition}
\begin{proof}
  This is the declared model definition of area In Square Meters.
\end{proof}
\begin{definition}[Frustum Figure Label]
  \label{def:physics:phyx-mini-0666:phyxminiproblems-problemphyxmini0666-frustumfigurelabel}
  \lean{PhyXMiniProblems.ProblemPhyXMini0666.FrustumFigureLabel}
  The three dimension labels printed in the supplied diagram
\end{definition}
\begin{proof}
  This is the declared model definition of Frustum Figure Label.
\end{proof}
\begin{definition}[Frustum Figure Feature]
  \label{def:physics:phyx-mini-0666:phyxminiproblems-problemphyxmini0666-frustumfigurefeature}
  \lean{PhyXMiniProblems.ProblemPhyXMini0666.FrustumFigureFeature}
  Named geometric features visible in the supplied frustum diagram
\end{definition}
\begin{proof}
  This is the declared model definition of Frustum Figure Feature.
\end{proof}
\begin{definition}[Supplied Frustum Figure]
  \label{def:physics:phyx-mini-0666:phyxminiproblems-problemphyxmini0666-suppliedfrustumfigure}
  \lean{PhyXMiniProblems.ProblemPhyXMini0666.SuppliedFrustumFigure}
  \uses{def:physics:phyx-mini-0666:phyxminiproblems-problemphyxmini0666-frustumfigurelabel, def:physics:phyx-mini-0666:phyxminiproblems-problemphyxmini0666-frustumfigurefeature}
  Presentation-level evidence transcribed from image 666.png. It contains no surface-area value and no answer-choice selection
\end{definition}
\begin{proof}
  This is the declared model definition of Supplied Frustum Figure.
\end{proof}
\begin{definition}[Right Circular Conical Frustum]
  \label{def:physics:phyx-mini-0666:phyxminiproblems-problemphyxmini0666-rightcircularconicalfrustum}
  \lean{PhyXMiniProblems.ProblemPhyXMini0666.RightCircularConicalFrustum}
  \uses{def:physics:phyx-mini-0666:phyxminiproblems-problemphyxmini0666-lengthquantity, def:physics:phyx-mini-0666:phyxminiproblems-problemphyxmini0666-suppliedfrustumfigure}
  The physical frustum and its independent derived observables. In particular, neither slantHeight nor curvedSurfaceArea is defined from the requested closed formula
\end{definition}
\begin{proof}
  This is the declared model definition of Right Circular Conical Frustum.
\end{proof}
\begin{definition}[Matches Supplied Frustum Figure]
  \label{def:physics:phyx-mini-0666:phyxminiproblems-problemphyxmini0666-matchessuppliedfrustumfigure}
  \lean{PhyXMiniProblems.ProblemPhyXMini0666.MatchesSuppliedFrustumFigure}
  \uses{def:physics:phyx-mini-0666:phyxminiproblems-problemphyxmini0666-frustumfigurelabel, def:physics:phyx-mini-0666:phyxminiproblems-problemphyxmini0666-frustumfigurefeature, def:physics:phyx-mini-0666:phyxminiproblems-problemphyxmini0666-suppliedfrustumfigure}
  Exact transcription of the labels and qualitative geometry in 666.png
\end{definition}
\begin{proof}
  This is the declared model definition of Matches Supplied Frustum Figure.
\end{proof}
\begin{definition}[Has Valid Frustum Parameters]
  \label{def:physics:phyx-mini-0666:phyxminiproblems-problemphyxmini0666-hasvalidfrustumparameters}
  \lean{PhyXMiniProblems.ProblemPhyXMini0666.HasValidFrustumParameters}
  \uses{def:physics:phyx-mini-0666:phyxminiproblems-problemphyxmini0666-lengthinmeters, def:physics:phyx-mini-0666:phyxminiproblems-problemphyxmini0666-rightcircularconicalfrustum}
  Positivity and radius ordering for the nondegenerate frustum in the image
\end{definition}
\begin{proof}
  This is the declared model definition of Has Valid Frustum Parameters.
\end{proof}
\begin{definition}[Satisfies Right Frustum Meridian Geometry]
  \label{def:physics:phyx-mini-0666:phyxminiproblems-problemphyxmini0666-satisfiesrightfrustummeridiangeometry}
  \lean{PhyXMiniProblems.ProblemPhyXMini0666.SatisfiesRightFrustumMeridianGeometry}
  \uses{def:physics:phyx-mini-0666:phyxminiproblems-problemphyxmini0666-lengthinmeters, def:physics:phyx-mini-0666:phyxminiproblems-problemphyxmini0666-rightcircularconicalfrustum}
  Pythagoras applied to a meridian cross-section of a right conical frustum. This law determines the slant height but does not mention the requested area
\end{definition}
\begin{proof}
  This is the declared model definition of Satisfies Right Frustum Meridian Geometry.
\end{proof}
\begin{definition}[Satisfies Frustum Curved Surface Area Law]
  \label{def:physics:phyx-mini-0666:phyxminiproblems-problemphyxmini0666-satisfiesfrustumcurvedsurfacearealaw}
  \lean{PhyXMiniProblems.ProblemPhyXMini0666.SatisfiesFrustumCurvedSurfaceAreaLaw}
  \uses{def:physics:phyx-mini-0666:phyxminiproblems-problemphyxmini0666-lengthinmeters, def:physics:phyx-mini-0666:phyxminiproblems-problemphyxmini0666-areainsquaremeters, def:physics:phyx-mini-0666:phyxminiproblems-problemphyxmini0666-rightcircularconicalfrustum}
  The general curved-surface-area law for a conical frustum, stated using the independent physical slant height rather than the answer's eliminated form
\end{definition}
\begin{proof}
  This is the declared model definition of Satisfies Frustum Curved Surface Area Law.
\end{proof}
\begin{definition}[Answer Choice]
  \label{def:physics:phyx-mini-0666:phyxminiproblems-problemphyxmini0666-answerchoice}
  \lean{PhyXMiniProblems.ProblemPhyXMini0666.AnswerChoice}
  Labels of the four answer choices in the source problem
\end{definition}
\begin{proof}
  This is the declared model definition of Answer Choice.
\end{proof}
\begin{definition}[answer Choice Coefficient]
  \label{def:physics:phyx-mini-0666:phyxminiproblems-problemphyxmini0666-answerchoicecoefficient}
  \lean{PhyXMiniProblems.ProblemPhyXMini0666.answerChoiceCoefficient}
  \uses{def:physics:phyx-mini-0666:phyxminiproblems-problemphyxmini0666-answerchoice}
  Dimensionless coefficient multiplying the common expression in each option
\end{definition}
\begin{proof}
  This is the declared model definition of answer Choice Coefficient.
\end{proof}
\begin{definition}[displayed Answer Value]
  \label{def:physics:phyx-mini-0666:phyxminiproblems-problemphyxmini0666-displayedanswervalue}
  \lean{PhyXMiniProblems.ProblemPhyXMini0666.displayedAnswerValue}
  \uses{def:physics:phyx-mini-0666:phyxminiproblems-problemphyxmini0666-lengthinmeters, def:physics:phyx-mini-0666:phyxminiproblems-problemphyxmini0666-rightcircularconicalfrustum, def:physics:phyx-mini-0666:phyxminiproblems-problemphyxmini0666-answerchoice, def:physics:phyx-mini-0666:phyxminiproblems-problemphyxmini0666-answerchoicecoefficient}
  Square-metre value of one displayed answer expression
\end{definition}
\begin{proof}
  This is the declared model definition of displayed Answer Value.
\end{proof}
% --- Archon declaration topology end ---
% --- Archon physics formalization source end ---
